\input{../tex-inputs/latex-leseansicht-vorspann}

\section[Arthur und Olga Schnitzler an Gustav Schwarzkopf, 30. 7. 1912]{L04531 Arthur und Olga Schnitzler an Gustav Schwarzkopf, 30. 7. 1912}
\nopagebreak\mylabel{L04531v}
\rehead{ }\normalsize\beginnumbering\briefempfaengerindex{Schwarzkopf, Gustav@\textsc{Schwarzkopf, Gustav}!zzzSchnitzler, Olga@\emph{von Olga Schnitzler}!1912-07-121@{12. 7. 1912}|(be}\briefempfaengerindex{Schwarzkopf, Gustav@\textsc{Schwarzkopf, Gustav}!zzzSchnitzler, Arthur@\emph{von Arthur Schnitzler}!1912-07-121@{12. 7. 1912}|(be}
\toendnotes[C]{\smallbreak\pagebreak[2]}
\correspDesc{Versand  durch Arthur Schnitzler, Olga Schnitzler am 12. 7. 1912 in Brijuni
\newline{}Übermittlung  am 12. 7. 1912 in Brijuni
\newline{}Erhalt  durch Gustav Schwarzkopf im Zeitraum [13. 7. 1912
                  – 17. 7. 1912?] in Wien}\toendnotes[C]{\smallbreak}
\Standort{DLA, A:Schnitzler, HS.1985.1.1897.}
\physDesc{Bildpostkarte, 204 Zeichen
\newline{}Handschrift Arthur Schnitzler: Bleistift, deutsche Kurrent
\newline{}Handschrift Olga Schnitzler: Bleistift, lateinische Kurrent
\newline{}Versand: Stempel: »\nobreak{}\oindex{Brijuni@\textbf{Brijuni}|pwk}Brioni, 30. 7. 12\nobreak{}«.  }\pstart{}{\pb}Herrn \textsc{Gustav Schwarzkopf}\pend{}\pstart{}\textsc{Wien}\oindex{Wien@\textbf{Wien}|pw}\pend{}\pstart{}\textsc{I. Tiefer Graben 17}\oindex{Wien@\textbf{Wien}!I., Innere Stadt@\textbf{I., Innere Stadt}!Tiefer Graben 17@\textbf{Tiefer Graben 17}|pw}\pend{}{\bigskip}
\pstart
           \noindent{}\centering{}{\pb}\textcolor{gray}{\textbf{Insel Brioni i.d. Adria\oindex{Brijuni@\textbf{Brijuni}|pw}}}\pend
           
\pstart
           \centering{}\textcolor{gray}{\textbf{Sonnenuntergang}}\pend
           \vspace{1em}
\pstart
           {\pb}30/7 12\pend
           \vspace{0.5em}
\pstart
           Sie ſind wohl ſchon wieder daheim? Wir fühlen uns hier ſehr wohl, grüßen Sie
               herzlichſt!\pend
           \pstart Ihr \spacefill\mbox{Arthur}\pend{}\selectlanguage{ngerman}\vspace{1em}
\pstart
           \noindent{}{[}hs. O. Schnitzler:{]} Hier müssen Sie einmal herkommen, es ist so schön!\pend
           \pstart \spacefill\mbox{Olga}\pend{}\selectlanguage{ngerman}\endnumbering\briefempfaengerindex{Schwarzkopf, Gustav@\textsc{Schwarzkopf, Gustav}!zzzSchnitzler, Olga@\emph{von Olga Schnitzler}!1912-07-121@{12. 7. 1912}|)be}\briefempfaengerindex{Schwarzkopf, Gustav@\textsc{Schwarzkopf, Gustav}!zzzSchnitzler, Arthur@\emph{von Arthur Schnitzler}!1912-07-121@{12. 7. 1912}|)be}\mylabel{L04531h}
\begin{anhang}
\end{anhang}\newcommand{\dateiname}{L04531}\newcommand{\titel}{Arthur und Olga Schnitzler an Gustav Schwarzkopf, 30. 7. 1912}\newcommand{\editorInnen}{Herausgegeben von Jahnke, SelmaMüller, Martin Anton}\input{../tex-inputs/latex-leseansicht-abspann}
